\chapter{Relations, Logic, and the Algebra of Records}
\label{chap:relations-logic-algebra}

The records are not just a list; they carry an algebra --- the operations ours to run, the count they preserve
the world's. Relations among entries compose and invert, channels
decompose beyond a single bit, a sheet of dependent values recoheres by recalculation, and a scratched disc
recovers its record by interleaving and parity. Five effects, all count-preserving, show the records closed
under their own operations --- relation, independence, channel, recompute, and correct.

\section{The Aristotle--De Morgan Effect: relations as an algebra}
\label{se:aristotle-demorgan}

Relational structure precedes symbolic formalism. Aristotle's logic was a logic of terms standing in relation
--- all men are mortal, Socrates is a man --- and De Morgan, two millennia on, saw that the relations
themselves form an algebra, with operations as definite as addition and multiplication. The binary relations
among the entries of a record --- which precedes which, which excludes which --- are not informal notes but
explicit algebraic objects, subject to composition, inversion, and duality. A relation $R$ has a converse,
$R^{-1}$, got by swapping each ordered pair,
\begin{equation}
  R^{-1} = \{\, (j, i) : (i, j) \in R \,\},
\end{equation}
and two relations compose, $R \circ S$, by chaining them through a shared middle,
\begin{equation}
  R \circ S = \{\, (i, k) : (i, j) \in R \text{ and } (j, k) \in S \text{ for some } j \,\}.
\end{equation}
Both operations preserve the count: neither adds nor removes an entry, only re-relates the ones already there
--- the composing and converting ours, the pairs they re-relate the world's.
And the duality of De Morgan holds throughout. For the logical connectives it is the familiar law,
\begin{equation}
  \neg(A \cup B) = \neg A \cap \neg B,
\end{equation}
the complement of a union the intersection of the complements; for relations it reappears as the rule that the
converse of a composition reverses the order of its factors,
\begin{equation}
  (R \circ S)^{-1} = S^{-1} \circ R^{-1},
\end{equation}
just as undoing two actions in turn undoes the last one first.

The honest floor is a finite count obligation --- composition and inversion are count-preserving operations on
the relations the record already holds --- the algebra ours to work, the entries already there the world's.
The reading is that relations form an algebra. The order and exclusion
among entries are not loose facts but a structure closed under converse and composition, an algebra the record
carries the moment it has more than one entry to relate.

\section{The Implied Orthogonality Effect: an independent axis, not a perpendicular}
\label{se:implied-orthogonality}

A record carries an information degree of freedom that is independent of its spacetime label, and the
independence must be stated with care. Project a reading onto where and when it was taken,
\begin{equation}
  \pi(\text{reading}) = (\text{here}, \text{now}),
\end{equation}
and the information content of the reading is a separate axis from that spacetime label: the same information
can sit at any here-and-now, and the same here-and-now can carry any information, so the two factor as an
independent product --- the projection ours to take, the independence it reveals the world's. A library makes the point plainly --- the content of a book is independent of the shelf it
stands on, the same text able to sit anywhere and any shelf able to hold any text --- and no one mistakes the
catalogue's independence of content from location for a claim that books are perpendicular to shelves. This is
what ``orthogonal to spacetime'' means here, and it is a set-theoretic independence, not a geometric
perpendicularity: there is no right angle, no extra spatial dimension, only the logical independence of two
coordinates.

The honest floor is a finite count obligation --- the information axis and the spacetime axis are independent
coordinates, their joint count the product of their separate counts --- the axes ours to name, the factoring
the world's. The reading is that information is an axis
of its own, independent of where and when. Strip the geometric connotation from ``orthogonal'': the content of
a reading is not perpendicular to spacetime in any spatial sense but independent of it, able to take any value
at any place and time --- a separate coordinate, not a separate direction.

\section{The Fessenden--Shannon Effect: a channel beyond binary}
\label{se:fessenden-shannon}

A channel need not speak in bits. Beyond the binary off/on distinction, a record admits a finite decomposition
into a larger alphabet of distinguishable symbols, and refining the alphabet --- adding symbols --- does not
inflate the count: the same record is re-expressed in more letters, each entry still a single symbol, the total
unchanged --- the alphabet ours to draw, the count it re-expresses the world's. On Christmas Eve of 1906, Reginald Fessenden put his own voice and a violin over the radio, where
there had been only the two-state clatter of Morse --- the first broadcast to carry not dots and dashes but a
continuum of sound resolved into symbols an ear could read. Shannon, four decades later, gave the capacity its
measure: the most a channel can carry is the count of distinguishable symbols it admits per use, no more,
whatever the symbols. Channel capacity is exactly this admissible decomposition --- not how fast a
fixed alphabet runs, but how many distinctions the channel can hold at once.

The honest floor is a finite count obligation --- a richer alphabet re-decomposes the record without changing
its count, and capacity is the count of admissible symbols --- the decomposition ours to pick out, the
distinctions the channel can hold the world's. The reading is that a channel is a decomposition,
not a fixed pair of states. The two-state bit is one alphabet among many; refining to a larger alphabet carries
more per symbol without adding to the count, and the capacity of a channel is the number of distinctions it can
admit at once.

\section{The Excel Effect: coherence by recalculation, not propagation}
\label{se:excel}

A record of dependent values recoheres by recalculation, not by a propagating wave. A spreadsheet is the
everyday model: a sheet of cells, each a formula in others, is not a collection of independent variables but a
dependency graph, and when one cell changes the rest do not update one after another like a row of toppling
dominoes. They update together. The program walks the dependency graph and recomputes every cell that depends,
directly or indirectly, on the one that changed, bringing the whole sheet into agreement in a single settling,
no cell sending its new value down a wire to the next, only a recalculation that re-establishes coherence
across the dependencies at once --- the change ours to make, the settling that answers it the world's. The
dependent effects are not signals; they are recalculations.

The honest floor is a finite count obligation --- the recalculation is additive under append, restoring global
coherence by a count, not a propagation --- the graph ours to set up, the agreement it settles into the
world's. The reading is that coherence is recomputed, not propagated. A change
does not send a wave through the record; it triggers a recalculation that settles the whole dependent structure
together, the global agreement restored by adding up the dependencies rather than passing a message along them.

\section{The Compact Disc Encoding: a record recovered across a scratch}
\label{se:compact-disc}

A record can be built to survive its own damage. A compact disc writes a finite alphabet of marks --- pits and
lands --- in an ordered sequence,
\begin{equation}
  e_1 \prec e_2 \prec \cdots,
\end{equation}
each position selecting one symbol, the word encoding the audio through a sequence of refinements governed by
cross-interleaved Reed--Solomon coding, the scheme called CIRC. Two ideas do the work. First, Reed--Solomon
parity: extra symbols are added to each block so that, knowing where the losses fell, a handful of missing
symbols can be solved back from the parity, the way two equations recover two unknowns. Second, interleaving:
before writing, the samples are shuffled so that neighbours in the music are scattered far apart on the disc.
Then a scratch --- which destroys a contiguous run of marks, hundreds in a row --- removes, once the samples
are de-interleaved back into order, only a scattered few from each block, few enough that parity recovers them.
A wound that would erase a long unbroken passage, written plainly, becomes a sprinkle of single losses, each
one repairable --- the shuffle ours, the scratch the world's. Coarsening the record preserves the count; the
correction restores it.

The honest floor is a finite count obligation --- the model claims the count survives coarsening and merge;
that a scratched record plays on is the engineering the section quotes, built on that surviving count --- the
scheme ours to build, the count that survives the world's. The reading is that error correction is a record
that protects its own
count. By spreading the order and carrying parity, the disc turns a contiguous scratch into a recoverable
scatter, and the record survives a wound that, written plainly, would have erased it.

\bigskip

With error correction, the claim is whole: the records form an algebra, closed under
the operations that act on them. Relations compose and invert, an information axis stands independent of
spacetime, a channel decomposes into any alphabet, a dependent sheet recoheres by recalculation, and a
scratched record recovers itself by parity --- and every one of these is count-preserving, an operation on the
entries the record already holds rather than an addition to them --- the operations ours, the count they never
touch the world's. Part VI turns next to computation itself ---
the universality, the limits, and the stability of the machine the records run on.
